\section{Boolean model}\label{sec:reduced_costs_boolean_model}
To derive reduced cost \eqref{eq:BP_redcost}, recall that the dual of the linear program (LP)
\begin{equation}\label{eq:LP_primal_general}
    \begin{split}
        \min & gx \\
        \st & Ax \geq b, \\
        & x \geq 0,
    \end{split}
\end{equation}
is given by
\begin{equation}\label{eq:LP_primal_dual}
    \begin{split}
        \max & yb \\
        \st & yA \leq g, \\
        & y \geq 0.
    \end{split}
\end{equation}
Denote by $A_k$ the $k$-th column of $A$.
The pricing problem asks whether there exists a column $k$ such that $yA_k > g$, which is equivalent to finding a $k$
such that reduced cost
\begin{equation}\label{eq:pricing_problem}
    \rho_k = g_k - yA_k < 0.
\end{equation}
Observe that LP relaxation of the master problem \eqref{eq:BM} can be written as
\begin{equation}\label{eq:BM_primal_vector}
    \begin{split}
        \min & \begin{bmatrix}
                   \bm{1} & d
        \end{bmatrix} \begin{bmatrix}
                          \zeta \\ w
        \end{bmatrix} \\
        \st & \begin{bmatrix}
                  I & A \\ 0 & -c
        \end{bmatrix} \begin{bmatrix}
                          \zeta \\ w
        \end{bmatrix}
        \geq \begin{bmatrix}
                 \bm{1} \\ -C
        \end{bmatrix}, \\
        & \zeta, w \geq 0,
    \end{split}
\end{equation}
where $\bm{1}$ is the all-one column vector,
\begin{align*}
    A_{ik} &= \Ind(k \in \K_i), & i \in \P, \; k \in \K, \\
    d_k &= \sum_{i \in \N} \Ind(k \in \K_i), & k \in \K. \\
\end{align*}
and $c$ is the row vector holding the complexity of each rule.
The dual is described by
\begin{equation}\label{eq:BM_dual_vector}
    \begin{split}
        \max & \begin{bmatrix}
                   \mu & \lambda
        \end{bmatrix} \begin{bmatrix}
                          \bm{1} \\ -C
        \end{bmatrix} \\
        \st & \begin{bmatrix}
                  \mu & \lambda
        \end{bmatrix} \begin{bmatrix}
                          I & A \\ 0 & -c
        \end{bmatrix}
        \leq \begin{bmatrix}
                 \bm{1} & d
        \end{bmatrix}, \\
        & \mu, \lambda \geq 0,
    \end{split}
\end{equation}
Therefore, the reduced cost is given by
\begin{equation}\label{eq:BP_redcost_vector_derivation}
    \begin{split}
        \rho_k &= d_k - \left(\mu A_k - \lambda c_k\right)
        \\ &= \sum_{i \in \N} \Ind(k \in \K_i) - \sum_{i \in \P} \mu_i\Ind(k \in \K_i) + \lambda c_k.
    \end{split}
\end{equation}

\section{Additive ramp loss model}\label{sec:reduced_costs_additive_ramp_loss}
The derivation of the reduced cost for the additive ramp loss model \eqref{eq:RM} is analogous to the derivation for the boolean model.
Consider \eqref{eq:BM_primal_vector}, only now $A$ is replaced by $\Gamma$, where
\begin{align*}
    \Gamma_{ik} &= \gamma_{ik}, & i \in \P, \; k \in \K,
\end{align*}
the $\bm{1}$-vector on the right-hand-side is replaced by a $\bm{2}$-vector,
and we choose
\begin{align*}
    d_k &= \sum_{i \in \N} \gamma_{ik}, & k \in \K,
\end{align*}
This formulation corresponds precisely to the additive ramp loss model \eqref{eq:RM}.
It follows that
\begin{equation}\label{eq:RP_redcost_vector_derivation}
    \rho_k = \sum_{i \in \N} \gamma_{ik} - \sum_{i \in \P} \mu_i\gamma_{ik} + \lambda c_k.
\end{equation}

\section{Non-additive ramp loss model}\label{sec:reduced_costs_nonadditive_ramp_loss}
The non-additive ramp-loss model formulation \eqref{eq:NRM} can be written as
\begin{equation}\label{eq:NRM_primal_vector}
    \begin{split}
        \min & \begin{bmatrix}
                   \bm{1} & d
        \end{bmatrix} \begin{bmatrix}
                          \zeta \\ w
        \end{bmatrix} \\
        \st & \begin{bmatrix}
                  \hat{I} & \hat{H} \\
                  I & \mathbf{2} \\
                  0 & -c
        \end{bmatrix} \begin{bmatrix}
                          \zeta \\ w
        \end{bmatrix}
        \geq \begin{bmatrix}
                 \bm{2} + \eta \\ \bm{2} \\  -C
        \end{bmatrix}, \\
        & \zeta, w \geq 0,
    \end{split}
\end{equation}
where $\hat{I} \in \R^{(\P \times \K)\times \P}$, $\hat{H}\in \R^{(\P \times \K) \times \K}$ such that
\begin{align*}
    \hat{I}_{[ik']i'} &= \Ind(i = i') & i, i' \in \P, \quad k' \in \K, \\
    \hat{H}_{[ik']k} &= \eta_{ik'} - \eta_{ik}, & i \in \P, \quad k', k \in \K,
\end{align*}
and $c, d$ are the same as before.
Furthermore, $\mathbf{2}$ is the all-two matrix on the left hand side
and the all-two column vector on the right hand side.
The dual then consists of variables $\mu \in \R_+^{\P \times \K}$, $\theta \in \R_+^{\P}$ and $\lambda \in \R_+$.
The reduced cost of a rule $k$ is given by
\begin{equation}\label{eq:NRP_redcost_vector_derivation}
    \begin{split}
        \rho_k &= d_k - \left(\mu \hat{H}_k + \theta\mathbf{2} - \lambda c_k\right) \\
        &= \sum_{i \in \N} \gamma_{ik} - \sum_{\substack{i \in \P, k' \in \K \\ \eta_{ik} < \eta_{ik'}}} \mu_{ik'}(\eta_{ik'}-\eta_{ik}) + \sum_{i \in \P} 2\theta_i + \lambda c_k. \\
        &= \sum_{i \in \N} \gamma_{ik} - \sum_{\substack{i \in \P, k' \in \K \\ \eta_{ik} < \eta_{ik'}}} \mu_{ik'}\eta_{ik'} + \sum_{\substack{i \in \P, k' \in \K \\ \eta_{ik} < \eta_{ik'}}} \mu_{ik'}\eta_{ik} + \sum_{i \in \P} 2\theta_i + \lambda c_k. \\
    \end{split}
\end{equation}
The chosen rule influences what elements are included in the sum.
This makes translating the reduced cost to a MIP non-trivial and complicated.


%%Observe that the second sum is a constant.
%%And the third sum is a linear combination of $\eta_{ik}$`s.
%\textcolor{red}{TODO: Problem: a new column also brings new rows, which include existing columns.
%And finding a column includes shadow prices of columns that don't exist. This will get complicated.}

The formulation \eqref{eq:NRMa} may be written as
\begin{equation}\label{eq:NRMa_primal_vector}
    \begin{split}
        \min & \begin{bmatrix}
                   \bm{1} & d & 0
        \end{bmatrix} \begin{bmatrix}
                          \zeta \\ w \\ u
        \end{bmatrix} \\
        \st &
        \begin{bmatrix}
            I & 0 & -\hat{H} \\
            0 & 0 & \hat{I}^u \\
            0 & 0 & -\hat{I}^u \\
            0 & \hat{I}^w & -I \\
            0 & -c & 0
        \end{bmatrix}
        \begin{bmatrix}
            \zeta \\ w \\ u
        \end{bmatrix}
        \geq \begin{bmatrix}
                 \bm{2} \\ 1 \\ -1 \\ 0 \\ -C
        \end{bmatrix}, \\
        & \zeta, w, u \geq 0,
    \end{split}
\end{equation}
where $\hat{H}, \hat{I}^u \in \R^{\P \times (\P \times \K)}$, $\hat{I}^w \in \R^{(\P \times \K) \times \K}$, given by
\begin{align*}
    \hat{H}_{i[i',k]} &= \eta_{ik} \Ind(i = i'), & i \in \P, \quad [i', k] \in \P \times \K, \\
    \hat{I}^u_{i[i'k]} &= \Ind(i = i'), & i \in \P, \quad [i', k] \in \P \times \K, \\
    \hat{I}^w_{[ik']k} &= \Ind(k = k'), & [i, k'] \in \P \times \K, \quad k \in \K,
\end{align*}
The dual then consists of variables
\begin{align*}
    \mu &\in \R_+^\P, \\
    \nu_1 &\in \R_+^\P, \\
    \nu_2 &\in \R_+^\P, \\
    \theta &\in \R_+^{\P \times \K}, \\
    \lambda &\in \R_+.
\end{align*}
The reduced cost of a rule $k$ is given by
\begin{equation}\label{eq:NRPa_redcost_vector_derivation}
    \begin{split}
        \rho_k &= d_k - \theta \hat{I}^w_k + \lambda c_k \\
        &= \sum_{i \in \N} \gamma_{ik} - \sum_{i \in \P} \theta_{ik} + \lambda c_k.
    \end{split}
\end{equation}

